\documentclass[10pt]{article}
\usepackage[margin=0.75in,top=0.6in,bottom=0.6in]{geometry}
\usepackage{amsmath,amssymb,booktabs,microtype,xcolor,url}
\usepackage[colorlinks=true,linkcolor=blue!60!black,urlcolor=blue!60!black,pdftitle={FRE-GY-9743 Assignment 5},pdfauthor={Raj Pawar}]{hyperref}
\newcommand{\code}[1]{\texttt{#1}}
\setlength{\parindent}{0pt}\setlength{\parskip}{4pt}\emergencystretch=3em\pagestyle{empty}
\makeatletter\renewcommand{\section}{\@startsection{section}{1}{\z@}{8pt plus 2pt}{3pt}{\normalfont\large\bfseries}}\makeatother
\begin{document}
{\Large\bfseries FRE-GY-9743 Assignment 5: Implicit Function Differentiation and a Valuation Engine}\\[1pt]
Raj Pawar, NYU Tandon School of Engineering, October 25, 2026.\\
Code: \href{https://github.com/pawarraj8888/FRE-GY-9743-Assignments-5/tree/Raj_Pawar_Branch}{github.com/pawarraj8888/FRE-GY-9743-Assignments-5}, branch \code{Raj\_Pawar\_Branch}.\\
Project page: \href{https://pawarraj8888.github.io/FRE-GY-9743-Assignments-5/}{pawarraj8888.github.io/FRE-GY-9743-Assignments-5}.
\vspace{2pt}\hrule

\section*{Part I: implicit function differentiation}
\textbf{Setting.} The model has an internal state $X_I\in\mathbb{R}^n$ (for a curve, the forward rates at the nodes; for a smile, the
interpolator's parameters). The market gives me $n$ quotes $X_M$, and for each quote the model has a pricing function $V_i(X_I)$.
Calibrating means finding the $X_I$ with $V(X_I)=X_M$: as many equations as unknowns.

\textbf{Q1: solving $V(X_I)=X_M$.} It is a square root-finding problem, $F(X_I):=V(X_I)-X_M=0$, and my default tool is Newton's method.
From a starting point $X_I^{(0)}$ (yesterday's calibration, or a flat state) I iterate
\[
J\bigl(X_I^{(k)}\bigr)\,\Delta^{(k)} = -F\bigl(X_I^{(k)}\bigr),\qquad X_I^{(k+1)}=X_I^{(k)}+\lambda\,\Delta^{(k)},\qquad
J=\frac{\partial V}{\partial X_I},
\]
solving the linear system rather than inverting $J$, with a step length $\lambda\le1$ (a line search, or Levenberg-Marquardt damping)
when I start far away, and bounds that keep the state admissible. I stop when $\|F\|$ is below a tolerance measured in quote units, a
small fraction of a basis point or of a vol point. The Jacobian comes analytically when the model exposes $\partial V/\partial X_I$,
which is what the gradient methods in this library are for; otherwise by finite differences, at $n$ extra pricings per iteration.
Two things make the problem easier in practice. If instrument $i$ only depends on the first $i$ state variables, $J$ is lower
triangular and Newton collapses into $n$ one-dimensional root searches done in order: that is bootstrapping a curve node by node, with
Brent or the secant method. And if each $V_i$ is monotone in the parameter it mostly depends on (a price in a vol, a swap rate in a
forward rate), the solution is unique and Newton converges in a handful of steps. Existence and uniqueness need $J$ non-singular at the
solution, which is also what the next question needs.

\textbf{Q2: $dX_I/dX_M$ at the solution.} The calibrated state is a function of the quotes, $X_I^*(X_M)$, defined implicitly by
$V\bigl(X_I^*(X_M)\bigr)=X_M$. Differentiating that identity with respect to $X_M$ and using the chain rule,
\[
\frac{\partial V}{\partial X_I}\,\frac{dX_I^*}{dX_M}=I \quad\Longrightarrow\quad \frac{dX_I^*}{dX_M}=J^{-1},\qquad
J=\frac{\partial V}{\partial X_I}\Big|_{X_I^*}.
\]
That is the implicit function theorem: no re-calibration and no bumping of quotes, just the inverse of the Jacobian I already had
from the last Newton step. Column $j$ of $J^{-1}$ says how every state variable moves when quote $j$ moves by one unit with the others
held fixed. It is also how model risk becomes market risk: for any trade with value $P(X_I)$,
\[
\frac{dP}{dX_M}=\frac{\partial P}{\partial X_I}\,J^{-1},
\]
so a vector of state sensitivities (what Part~II computes) becomes hedge ratios in the quoted instruments after one linear solve.
In this assignment the state \emph{is} the quote (the IFR nodes are the inputs), so $J=I$ and the risk report needs no Jacobian; a
curve calibrated to swap rates would need the $J^{-1}$ step. If $J$ is nearly singular the entries of $J^{-1}$ blow up, and the theorem
is telling me that the calibration itself is unstable, not that the hedges are large.

\textbf{Q3: three OTM options and a vol interpolator.} The quotes are the undiscounted prices $P_1,P_2,P_3$ at
$K_1<K_{\mathrm{atm}}<K_3$ (an OTM put, the ATM option, an OTM call), the state is $\theta\in\mathbb{R}^3$, and the model price of
strike $K_i$ is $V_i(\theta)=\mathrm{Black}\bigl(F,K_i,T,\sigma(K_i;\theta)\bigr)$ without the discount factor. Three prices and three
parameters, so the recipe from Q1 applies, with the chain rule for the Jacobian:
\[
J_{ij}=\frac{\partial V_i}{\partial\theta_j}=
\underbrace{\frac{\partial\,\mathrm{Black}_i}{\partial\sigma}}_{\text{vega at }K_i}\;\frac{\partial\sigma(K_i;\theta)}{\partial\theta_j}.
\]
I would not actually run Newton on prices, though. Black's price is strictly increasing in $\sigma$, so each $P_i$ inverts to a unique
implied vol $\sigma_i^{\mathrm{mkt}}$ with a one-dimensional solve, and the calibration becomes $\sigma(K_i;\theta)=\sigma_i^{\mathrm{mkt}}$
for $i=1,2,3$: the interpolator has to pass through three points. If the interpolator is linear in its parameters, say a quadratic in
log-moneyness $x=\ln(K/F)$ with $\sigma(K;\theta)=\theta_0+\theta_1x+\theta_2x^2$, this is a $3\times3$ linear system with the constant
matrix $[\,1,\ x_i,\ x_i^2\,]$, solved exactly in one step, and $\theta$ is unique whenever the three strikes are distinct. For a nonlinear
family such as SABR with $\beta$ fixed, $\theta=(\alpha,\rho,\nu)$, Newton takes a few iterations from a start read off the quotes
($\alpha$ from the ATM level, $\rho$ from the slope, $\nu$ from the curvature), with $|\rho|<1$ and $\nu>0$ enforced.

How $\sigma(K)$ moves with $\theta$ is the gradient of the interpolator, $\partial\sigma(K;\theta)/\partial\theta$, one row of three
numbers per strike; for the quadratic it is just the basis $(1,x,x^2)$. The more useful question is how the whole smile moves when a quote
moves, and that is Q2 again: $d\sigma(K)/dP=\bigl(\partial\sigma(K)/\partial\theta\bigr)J^{-1}$, or in vol terms
$d\sigma(K)/d\sigma^{\mathrm{mkt}}=\bigl(\partial\sigma(K)/\partial\theta\bigr)\bigl[\partial\sigma(K_i;\theta)/\partial\theta\bigr]^{-1}$,
three weights that say how much of each quoted option reproduces a vol move at $K$ (for an interpolator that can represent a flat smile
they sum to one, because shifting the three quotes together shifts the whole smile).

\textbf{What $\theta$ means to a trader.} Nobody quotes $\theta_0,\theta_1,\theta_2$. The desk quotes the ATM vol, the risk reversal
$\sigma(K_3)-\sigma(K_1)$ and the butterfly $\tfrac12\bigl(\sigma(K_1)+\sigma(K_3)\bigr)-\sigma(K_{\mathrm{atm}})$: level, slope and
curvature of the smile. With three parameters and three quotes the map between them is one-to-one, and for symmetric strikes
$x_3=-x_1=d$, $x_{\mathrm{atm}}=0$ it is explicit: $\mathrm{ATM}=\theta_0$, $\mathrm{RR}=2\theta_1d$, $\mathrm{BF}=\theta_2d^2$. So
intuitively $\theta_0$ is the level, $\theta_1$ the skew (the tilt of the smile, the spot-vol correlation, $\rho$ in SABR) and $\theta_2$
the smile (its convexity, the vol of vol, $\nu$ in SABR). If I parameterise the interpolator directly in (ATM, RR, BF), then
$\partial\sigma(K)/\partial\theta$ reads as ``vol points at $K$ per vol point of ATM, of RR and of BF'', which is the form in which the
trader hedges it with the three quoted options.

\section*{Part II: valuation engine and first-order risk}
I completed the five \code{TODO}s. In \path{fixedincomelib/yield_curve/yield_curve_model.py}: (1)~the component's local gradient. With
$D_{\mathrm{local}}(T)=\exp\bigl(-\int_0^T f\bigr)$ and $f$ piecewise constant, $\partial D/\partial f_i=-D\,w_i$, where $w_i$ is the
length of node $i$'s bucket inside $[0,T]$ and is exactly what \code{gradient\_of\_integrated\_value\_wrt\_ordinate(0,T)} returns; the
method writes $-\mathrm{scaler}\cdot D\,w$ into the caller's array, replacing it or adding to it. (2)~The model's gradient walks the
reference chain the way \code{discount\_factor} does. Since $D_{\mathrm{funding}}=D_{\mathrm{SOFR}}\,D_{\mathrm{spread}}$, the product rule
gives the SOFR block $D_{\mathrm{spread}}\,\partial D_{\mathrm{SOFR}}/\partial f$ and the funding block
$D_{\mathrm{SOFR}}\,\partial D_{\mathrm{spread}}/\partial s$; I recurse with the running product carried in the scaler, and when overwriting
I zero every block first so that a component off the chain comes back as zeros rather than stale numbers. In
\path{fixedincomelib/yield_curve/valuation_engine.py}: (3)~the constructor reads the \code{FUNDING INDEX PARAMETER}, resolves the
funding identifier for the product's currency and stores $C=N\alpha$; (4)~\code{calculate\_value} sets $PV=C\,D(T_p)$ and cash $0$ for a
future payment, $PV=\mathrm{cash}=C$ on the value date, and zeros for a past payment; (5)~\code{calculate\_risk} is $C$ times the model
gradient, with $C$ passed through as the scaler, and zero blocks (or nothing added) when there is no discount factor to differentiate.

The notebook's case: value date 2026-01-02, payment 2026-10-02, so $T=273/365=0.747945$.
The SOFR IFR nodes sit at $0.246575$, $1$ and $2$ years and the spread nodes at $0.495890$ and $1$ year, so the bucket lengths inside
$[0,T]$ are $(0.246575,\ 0.501370,\ 0)$ and $(0.495890,\ 0.252055)$: the third SOFR node carries no risk, and the first spread node does
even though its spread is zero. With $D_{\mathrm{SOFR}}=0.972921$, $D_{\mathrm{spread}}=0.999748$,
$D_{\mathrm{funding}}=0.972676$ and $C=1{,}000{,}000\times91/360=252,777.78$, the cashflow is worth
$PV=245,870.91$ with no cash today. The last column is the signed PV change for a one basis point rise in the node, the
\code{UNIT} scaling the risk report applies on display.

\begin{center}\small\setlength{\tabcolsep}{5pt}
\begin{tabular}{@{}llrrrrrr@{}}\toprule
Component & Node & IFR & $w_i$ & $\partial D_{\mathrm{local}}/\partial f_i$ & $\partial D_{\mathrm{funding}}/\partial f_i$ & $\partial PV/\partial f_i$ & per bp\\ \midrule
SOFR-1B & 2026-04-02 & 0.030 & 0.246575 & -0.239898 & -0.239838 & -60,625.70 & -6.0626 \\
SOFR-1B & 2027-01-02 & 0.040 & 0.501370 & -0.487793 & -0.487670 & -123,272.26 & -12.3272 \\
SOFR-1B & 2028-01-02 & 0.050 & 0.000000 & 0.000000 & 0.000000 & 0.00 & 0.0000 \\
SOFR-1B-FLAT & 2026-07-02 & 0.000 & 0.495890 & -0.495765 & -0.482341 & -121,925.03 & -12.1925 \\
SOFR-1B-FLAT & 2027-01-02 & 0.001 & 0.252055 & -0.251991 & -0.245168 & -61,972.94 & -6.1973 \\
\bottomrule\end{tabular}\end{center}

The notebook runs top to bottom with the professor's cells untouched; its printed gradients, PV and risk table agree with the hand
calculation above to the six printed decimals. \path{tests/test_yield_curve_gradient.py} (16 tests) checks the
local and model gradients against the bucket-length formula and against central bump-and-reval through
\code{perturb\_model\_parameter}, the product rule, the overwrite and accumulate semantics, block allocation, serialization and the
validation errors; \path{tests/test_valuation_engine.py} (12 tests) checks the engine's PV, cash and risk for future,
same-day and past payments, the sign of a short position, the registry dispatch and the PV and risk reports.
\end{document}
